\section{Polynomial scale profiles}\label{sec:profiles}

We work along the contradiction sequence, with
\[
 N=|P|,\qquad B=N^{1/d},\qquad A=B^2,\qquad M=o(A).
\]
In this section a further parameter $\delta\to0$ is given, and
$E=\delta N/A$.  The parameter will be the total directional degree in
the sparse-cones argument.  All constants below may depend on $d$.
Our aim is to partition the sample into pieces on which successive
polynomial cutting degrees can be recorded uniformly.  The resulting
profiles will bound component masses from above and polynomial
restriction ranks from below.

\subsection{Conventions and the partition}

For a width $s\to\infty$, $s\le\log B$, put
\begin{equation}\label{prof:testing-degree}
 T_s(b)=B\exp(bs).
\end{equation}
Degrees are rounded to integers.  Throughout, an estimate
$U\ge V\exp(-o(s))$ means
\[
 \liminf\frac{\log(U/V)}s\ge0.
\]
Equivalently, every fixed $\varepsilon>0$ can replace $o(s)$ by
$\varepsilon s$, for all sufficiently large indices.  The analogous
convention applies to upper bounds.  We only use such errors in the
presence of a subsequently specified strict positive margin.

A \emph{piece} below is a subset of the sample, together with one path
of polynomial cuts.  The entire path, rather than a different path at
each point, is attached to the piece.  At every point on the piece each
successive cut is proper on every component of the preceding local
system.  We refer to the irreducible components of that system meeting
the piece as its \emph{carriers}.  Different points can have different
carriers.  Properness at the counted points will suffice for every
degree estimate.

At testing degree $D$, the normalized time is $\log(D/B)/s$.
We record each cut at the testing time when it is made; its actual
degree may be smaller.  When the $d$ recorded times on a path have
subsequential limits $b_1,\ldots,b_d$,
we call each open interval between consecutive distinct limits a
\emph{plateau}.  Its number of completed cuts, and hence its local
dimension, is constant.  We also use initial and final plateaux for
the intervals before the first and after the last limiting cut.

\begin{proposition}[Multiscale partition]\label{prof:partition}
There are $s_*\to\infty$ and partitions of all but $o(N)$ sample points
into a core and non-core pieces with the following properties.
\begin{enumerate}
\item The core is simultaneously isolated by equations of a degree $L$
such that $EL^2=o(N)$.
\item Each non-core piece has a dyadic width
$s=2^{-a}\log B\ge s_*$ and size $n=N\exp(-o(s))$.  Its path has $d$
proper hypersurface cuts.  Along any sequence of pieces, after passing
to a subsequence, their normalized cut times converge to
\[
 -1\le b_1\le\cdots\le b_d\le1,\qquad
 \max_i|b_i|\ge\tfrac12.
\]
The degree of cut $i$ is at most
$D_i\exp(o(s))$, where $D_i=B\exp(b_i s)$.
\item Let $b$ lie strictly between consecutive distinct limiting cut
times, and let $j=\#\{i:b_i<b\}$ and $l=d-j$.  At each point of the
piece the current system is pure of local dimension $l$.  Uniformly
over polynomials of degree $O(T_s(b))$, the number of points where such
a polynomial vanishes but contains none of the local carriers is
$o(n)$.  The same assertion holds before the first cut whenever
$T_s(b)\ge1$ and $b<b_1$, and it is automatic after the last cut.
\end{enumerate}
The assertions about error terms and negligible proportions hold along
every sequence of pieces.  In particular, a violating choice of one
piece at each index is not permitted.
\end{proposition}

\begin{proof}
Choose $s_*\to\infty$ sufficiently slowly that
\[
 s_*\le\tfrac14\log B,\qquad \delta\exp(6s_*)\longrightarrow0.
\]
Start with width $\log B$, halve the width at each step, and stop on
entry to the first width in $[s_*,2s_*]$.  On entry to a step of width
$s$, the unresolved set is isolated by equations of degree
\begin{equation}\label{prof:entry-upper}
 B\exp((1+o(1))s).
\end{equation}
This is true initially by Lemma~\ref{poly:separators}, since
$M=o(B^2)$.  Except at the initial step, the unresolved set also has no
nonzero polynomial of degree at most $B\exp(-s)$ capturing as many as
$N\exp(-\sqrt{2s})$ points.

At a step which is not the final one, first remove greedily batches of
at least $N\exp(-\sqrt{s})$ points on zero sets of nonzero polynomials
of degree at most $B\exp(-s/2)$.  These batches are sent to the
regularization below.  From the remaining points remove greedily
batches of at least $N\exp(-\sqrt{s})$ points which occur as isolated
zero-dimensional components of some system of degree at most
$B\exp(s/2)$.  Only the union of these latter batches proceeds to the
next window.  The final remainder is sent to regularization if its
size is at least $N\exp(-s^{1/3})$, and otherwise is discarded.

There are at most $\exp(\sqrt{s})$ batches of either kind.  Isolated
point interpolation in Lemma~\ref{poly:interpolation} gives equations of
degree $O(B\exp(s/2))$ for precisely each batch of the second kind.
Indeed, from normalized separators $q_i$ for its points $p_i$ use
$1-\sum_iq_i$ and all coordinates of $(x-p_i)q_i$.  Their common zero
set is precisely the batch.  A union of finitely many such zero sets
is cut out by products, taking one equation from each system; the
degrees add.  Thus the union proceeding to the next step has isolation
degree at most
\[
 B\exp(s/2+O(\sqrt{s})).
\]
This is~\eqref{prof:entry-upper} at width $s/2$.  The first greedy
operation supplies exactly the asserted lower boundary there.
Consequently the final core has isolation degree at most
$B\exp(3s_*)$ for large indices, and hence $EL^2=o(N)$.

The shrinking windows have produced the core and the batches that
will become non-core pieces.  It remains to give each retained batch
one path of proper cuts and to ensure that a proper polynomial test
between cut times captures a vanishing fraction of the piece.  We
now enforce this condition while discarding only $o(N)$ points.
For each width $s$ use a grid in normalized time from $-1$ to
$1+o(1)$, containing the upper isolation degree, with mesh tending to
zero and $G=O(1+\sqrt{\log s})$ grid points.  Index the potential steps
by increasing grid time, and, within one grid time, by decreasing
positive dimension.  There are $K=O_d(G)$ indices.  Choose thresholds
$\lambda_i=\exp(-a_i)$ with
\begin{equation}\label{prof:lambda-hierarchy}
 a_K=s^{1/10},\qquad a_i=10a_{i+1}.
\end{equation}
The same indexed thresholds are used in all trees of this width.

In a node of dimension $k$ and original size $m$, repeatedly extract
at least $\lambda_i m$ remaining points where a polynomial of the
current testing degree vanishes properly on every local component.
An extracted set becomes a child, with that polynomial imposed, and
is processed in dimension $k-1$ at the same time.  Here $m$ stays the
original mass of this node during all its extractions.  The
unextracted residual advances to the next time, unless its size is
less than $\sqrt{\lambda_i}m$, in which case it is discarded.  At the
last time, at most $d$ generic combinations of the isolation equations
finish the path: on any positive-dimensional local component through
a counted point, some isolation equation is nonidentically zero, so
a generic combination is a proper cut simultaneously on all such
components.

Every path visits an index at most once.  Extraction decreases
dimension; taking a residual increases time.  Sibling extractions
are alternative branches, not repeated steps on one path.  Therefore
\[
 \sum_i a_i=s^{1/10+o(1)}=o(s^{1/3}),\qquad
 \sum_{r>i}a_r<\frac{a_i}{9}.
\]
If a residual survives at step $i$, every later terminal piece inside
it has size at least
$m\exp(-a_i/2-\sum_{r>i}a_r)$.  A proper cut available when
extraction stopped captured fewer than $m\exp(-a_i)$ residual points.
Its captured fraction in that terminal piece is consequently at most
\begin{equation}\label{prof:fragmentation-bound}
 \exp\left(-\frac{a_i}{2}+\sum_{r>i}a_r\right)
 \le \exp(-7a_i/18)\longrightarrow0.
\end{equation}
This proves the required no-capture assertion on any limiting open
plateau.  To allow a fixed multiplicative factor in the testing
degree, choose a slightly later grid time which is still inside that
plateau.

At a fixed index the nodes in all these trees are disjoint.  Their
discarded residual mass is at most $\exp(-a_i/2)$ times their total
mass.  Thus the regularization discards at most
\begin{equation}\label{prof:window-discard}
 C_d N(1+\sqrt{\log s})\exp(-s^{1/10}/2)
\end{equation}
points at width $s$.  This estimate is summable over windows even if
$s_*$ grows arbitrarily slowly.  To see this, reverse their order and
bound by widths $2^m u$, $m\ge0$, with $u\ge s_*$.  Put $t=u^{1/10}$
and $c=(\log2)/10$.  The inequalities
\[
 2^{m/10}\ge1+cm,\qquad
 \sqrt{\log(2^m u)}\le\sqrt{\log u}+\sqrt{m\log2}
\]
bound the sum of the fractions in~\eqref{prof:window-discard} by
\[
 C_d e^{-t/2}\sum_{m\ge0}
 (1+\sqrt{\log u}+\sqrt{m\log2})e^{-ctm/2}
 =O_d((1+\sqrt{\log u})e^{-u^{1/10}/2})=o(1).
\]
The separately discarded small remainders have fractions
$\exp(-s^{1/3})$ and are summable by the same argument.  No factor
equal to the number of windows is used.

The discarded mass is therefore negligible across all windows.
We check the remaining size and cut-time conditions.
Each retained terminal piece has size at least its initial batch size
times $\exp(-o(s^{1/3}))$, and hence is $N\exp(-o(s))$.  The inherited
lower boundary is negligible even relative to the smallest such
piece, since
\[
 \frac{N\exp(-\sqrt{2s})}
 {N\exp(-\sqrt{s}-o(s^{1/3}))}
 =\exp(-(\sqrt2-1)\sqrt{s}+o(s^{1/3}))\to0.
\]
There are exactly $d$ cuts on a path.  Their times have subsequential
limits $b_1\le\cdots\le b_d$ in $[-1,1]$.  For a first-type batch,
some cut occurs by limiting time $-1/2$: otherwise its defining
polynomial would capture the entire piece in an initial
full-dimensional no-capture interval.  For a remainder piece the
last cut cannot have limiting time less than $1/2$.  Otherwise its
path would isolate more than $N\exp(-\sqrt{s})$ points by equations
of degree at most $B\exp(s/2)$, contradicting termination of the
second greedy operation.  This proves the imbalance and all stated
properties.  The estimates were uniform over nodes and pieces, which
also proves the final sequence quantifier.
\end{proof}

\subsection{Degree budgets and propagation of rank}

The partition supplies paths and intervals on which proper tests
capture few points.  We now prove that their degree scales are
balanced and control both component masses and restriction ranks.
In particular, small exceptional component degree will imply small
sample mass.

For a limiting profile define
\begin{equation}\label{prof:profile-functions}
 D_i=B e^{b_i s},\qquad e_j=\prod_{i\le j}D_i,\qquad e_0=1,
 \qquad \sigma(b)=\prod_{i=1}^d\min(T_s(b),D_i).
\end{equation}
On a plateau after $j$ cuts, $\sigma(b)=e_jT_s(b)^{d-j}$.
Figure~\ref{fig:scale-profile} illustrates this relation.

\begin{proposition}[Budgets and balanced profiles]\label{prof:budget}
For every non-core profile,
\begin{equation}\label{prof:balance}
 \sum_{i=1}^d b_i=0,\qquad b_1<0<b_d,\qquad b_1+b_2<0.
\end{equation}
On a plateau after $j$ cuts, its carrier family has degree at most
$e_j\exp(o(s))$.  Every subcollection of total degree $a$ meets the
piece in at most
\begin{equation}\label{prof:mass-budget}
 a\left(\prod_{i>j}D_i\right)\exp(o(s))
\end{equation}
points.  This counts all points touching those components, whether
or not they have been assigned to them.  Finally, for every fixed
$\alpha>0$, every subset $S$ of at least $\alpha n$ piece points has
\begin{equation}\label{prof:rank-budget}
 H_S(T_s(b))\ge\sigma(b)\exp(-o(s))
\end{equation}
at each interior testing time, and with the usual limiting
interpretation at a cut time.
\end{proposition}

\begin{proof}
Apply B\'ezout successively, retaining the proper intersections
through counted points.  It gives the carrier degree bound and,
after applying the later cuts to a specified carrier subcollection,
\eqref{prof:mass-budget}.  Applied to the entire path it gives
$n\le(\prod_iD_i)\exp(o(s))$.  Since $n=N\exp(-o(s))$ and $N=B^d$,
this proves $\sum_i b_i\ge0$.

We prove~\eqref{prof:rank-budget} by propagation through the finitely
many plateaux; the assertion is always for every fixed positive
density subset, including subsets chosen during the proof.  Before
the first cut, no nonzero polynomial of the testing degree can
vanish on $S$.  Its envelope is therefore the ambient space, whose
rank is comparable to $T_s(b)^d$.

On a plateau of dimension $l$, let $V$ be the testing-degree envelope
of $S$. At all but $o(n)$ points of $S$, at least one local carrier
is contained in $V$. Indeed, if no carrier through a point is
contained, a generic linear combination of the envelope equations
is proper on all those carriers and vanishes at the point. There
are finitely many components and sample points, so the combination
can be chosen simultaneously; no-capture then applies.

For this envelope, choose at each such point one contained local
carrier, and choose one irreducible component of $V$ containing each
distinct carrier so chosen. Let the points inherit the component
chosen for their carrier. Group all irreducible components of $V$ by
dimension, counting the points by these choices; components receiving
no points remain in the degree of their group. Some dimension $h\ge l$
group, of total degree $a$, covers a fixed positive fraction of $S$.
Writing $S'$ for that fraction and $T=T_s(b)$, the component Hilbert
bounds give
\[
 H_S(T)=H_V(T)\gtrsim aT^h,\qquad
 H_{S'}(T')\lesssim a(T'+1)^h\quad(T'<T).
\]
The lower estimate uses that the group consists of components of
the envelope at its defining degree $T$; the upper estimate uses
only that $S'$ lies on that group.  Eliminating $a$ gives
\[
 H_S(T)\gtrsim H_{S'}(T')
       \left(\frac{T}{T'+1}\right)^h.
\]
Take $T'=T_s(b')$ at a time $b'$ just before the cut starting the
current plateau, and apply~\eqref{prof:rank-budget} there to $S'$.
The displayed inequality increases that lower bound by a factor
at least $(T/T_s(b'))^h$, to exponential accuracy.  Since $h\ge l$,
this is at least the growth required by $\sigma$ on the plateau.
Take $b'$ arbitrarily close to the
boundary; the logarithmic endpoint error can be made smaller than
any prescribed fixed multiple of $s$.  This proves the next
plateau, including after a group of simultaneous limiting cuts.

There is one possible boundary convention: if the first cut has
testing degree only $\exp(o(s))$, there may be no preceding testing
time of degree at least one.  Rank one at that boundary is already
the required bound to exponential accuracy, even if several first
cuts have that same time.  Use it as the start of the induction.

The same envelope argument advances across the final cut, with
required growth exponent zero.  Since rank is at most $N$, it gives
$\prod_iD_i\le N\exp(o(s))$, proving the reverse inequality
$\sum_i b_i\le0$.  The profile is nonzero by the imbalance in
Proposition~\ref{prof:partition}.  Its monotonicity and zero sum give
$b_1<0<b_d$.  If $b_1+b_2\ge0$, then $b_2>0$ and the $d-2$ later
entries are at least $b_2$, contradicting the zero sum because
$d>2$.  This proves~\eqref{prof:balance}.
\end{proof}

\begin{figure}[htbp]
\centering
\begin{tikzpicture}[x=5.7cm,y=.65cm,font=\small,
  profile/.style={very thick,draw=linkblue},
  cut/.style={densely dashed,draw=linkblue!55}]
 \fill[figureblue] (-.4,0) rectangle (.3,4.25);
 \draw[-{Latex[length=2mm]}] (-1.07,0) -- (-1.07,4.6)
   node[above right] {Carrier dimension $l=d-j$};
 \foreach \l in {0,1,2,3,4}{
   \draw (-1.09,\l) -- (-1.05,\l);
   \node[left] at (-1.09,\l) {$\l$};
 }
 \foreach \b/\i in {-.8/1,-.4/2,.3/3,.9/4}{
   \draw[cut] (\b,-.65) -- (\b,4.25);
   \node[below] at (\b,-.65) {$b_\i$};
 }
 \draw[-{Latex[length=2mm]}] (-1.07,-.65) -- (1.14,-.65)
   node[below left,yshift=-13pt] {Normalized time $b$};
 \node[below] at (-1,-.65) {$-1$};
 \node[below] at (1,-.65) {$1$};
 \draw[profile] (-1,4) -- (-.8,4);
 \draw[profile] (-.8,3) -- (-.4,3);
 \draw[profile] (-.4,2) -- (.3,2);
 \draw[profile] (.3,1) -- (.9,1);
 \draw[profile] (.9,0) -- (1,0);
 \draw[profile,dotted] (-.8,4) -- (-.8,3);
 \draw[profile,dotted] (-.4,3) -- (-.4,2);
 \draw[profile,dotted] (.3,2) -- (.3,1);
 \draw[profile,dotted] (.9,1) -- (.9,0);
 \node[align=center] at (-.05,3.35)
   {$(j,l)=(2,2)$\\[2pt]$\sigma(b)=e_2T_s(b)^2$};
\end{tikzpicture}

\caption{An illustrative limiting profile in dimension four, with fixed
window width $s$ and cut times $(b_1,b_2,b_3,b_4)=(-0.8,-0.4,0.3,0.9)$.
The degree scales are $D_i=T_s(b_i)$, where $T_s(b)=B e^{bs}$, and
$e_j=\prod_{i\le j}D_i$. On the highlighted open plateau $(b_2,b_3)$,
two cuts have occurred and $\sigma(b)=e_2T_s(b)^2$.
The staircase shows carrier dimension; its vertical drops are schematic.
General profiles may have coincident limiting cuts.}
\label{fig:scale-profile}
\end{figure}
\FloatBarrier

\subsection{Testing individual components and products}

\begin{proposition}[Regularity and individual cuts]\label{prof:regularity}
Inside an open plateau, individual proper cuts of degree $O(T_s(b))$
on the different carriers collectively capture only $o(n)$ points.
At a slightly later interior testing time there are common equations
for the entire carrier family whose Jacobian has full codimension
rank at all but $o(n)$ points.  At those surviving points the carrier
is unique and smooth, and its tangent space is testable by equations
and minors of the same interior degree scale.

The assertions also hold for the union of two such families, and
for unions from $\exp(o(s))$ pieces with the same limiting profile
and width, provided each piece retains its own assigned masses.
Common equations vanish on the whole union.  In particular, two
different members of that union cannot meet at a surviving point
where the common Jacobian has full rank.
Flat carriers which are not defined over $\R$ may be discarded
with negligible mass.
\end{proposition}

\begin{proof}
Every carrier used here is an irreducible component of its fixed
path locus.  Indeed, at every point of the piece on that carrier,
the successive proper cuts give the current local dimension.  A
larger irreducible component containing the carrier would pass
through that point and contradict this local dimension.  The path
locus may have larger components elsewhere; they remain in its
equations.

Let the plateau start at time $\beta<b$, and choose
$b<b'<b''$ in the same plateau.  The carrier components are therefore
isolated components of their cut system, whose equations have degree
at most $B\exp((\beta+o(1))s)$.  Lemma~\ref{poly:pooling}, used at degree
$T_s(b')$, pools proper component cuts of degree $O(T_s(b))$ into
one polynomial.  It is nonidentically zero except on components of
total degree at most
\[
 e_j\exp(-(b'-b)s+o(s)).
\]
By~\eqref{prof:mass-budget} and $\prod_iD_i=N$, all points of the
piece touching those exceptional components have mass $o(n)$.
At every remaining counted point the pooled polynomial is proper
on every local carrier.  No-capture at $b'$ proves the assertion.

Apply Lemma~\ref{poly:regularity} at degree $T_s(b')$.  Since the
cut-system equations have degree at most $B\exp((\beta+o(1))s)$,
the components on which the common equations fail to have generic
full rank have total degree at most
\[
 e_j\exp(-(b'-\beta)s+o(s)).
\]
The mass budget removes all piece points touching those components
at a cost of $o(n)$.  On the other components a generic linear
combination of maximal minors is nonidentically zero.  Its zero set
contains every rank failure, and its degree is $O(T_s(b'))$.
The preceding individual-cut conclusion applies to this linear
combination using the remaining interior margin before $b''$, and
removes its zeros at another $o(n)$ points.  At a surviving point the
implicit function theorem gives a smooth locus of the correct dimension; any two
carrier components through that point would have the same germ
and hence would coincide.  The equations vanish on the whole
family, including components whose generic regularity failed.

For pooled pieces, use the union-of-systems versions of
Lemmas~\ref{poly:pooling} and~\ref{poly:regularity}.  Each carrier
remains a component of its own path system.  On a generic transverse
fiber, Lemma~\ref{poly:union-interpolation} interpolates the selected
isolated lists, including first jets, with an effective bound equal
to the sum of the individual system bounds: separators for the other
lists are multiplied, using their exact-isolation equations for a
test point outside a list.  If these bounds are $L_\nu$, then for
$\exp(o(s))$ pieces with the common limiting profile,
\[
 L_*=\sum_\nu L_\nu\le B\exp((\beta+o(1))s).
\]
The strict interior gaps above absorb this change in the effective
bound and the $\exp(o(s))$ changes in the degree budgets.

Mass accounting for such a union is performed separately in each
original piece.  Its own later cuts bound all of its points
touching exceptional components of its own current system.  We
do not apply one piece's later cuts to points assigned to another
piece.  Rank failures on every good assigned component are then
discarded by that piece's no-capture assertion.  The equations
still vanish on the whole union, so full rank at a surviving point
also excludes intersections with a carrier from another piece.
The sequence quantifier makes the discarded fractions uniform
over these pieces; summing their masses creates no loss equal to
the number of pieces.

Finally, the real points of a nonreal complex affine flat lie in
a proper complex affine linear section of that flat.  Choose
one such degree-one proper cut for each nonreal flat, and no
cut on the other carriers.  The individual-cut assertion just
proved makes their total assigned mass negligible.
\end{proof}

For the next lemma, write $S_x,S_y$ for the two point sets and
$n_x=|S_x|$, $n_y=|S_y|$.  A dense graph is a subset
$\Gamma\subset S_x\times S_y$ of size at least $\alpha n_xn_y$
for some fixed $\alpha>0$.  A label map is a rational map $\phi$
from pairs to a fixed affine space, with bounded-degree coordinate
numerators and denominators on the chosen charts.  The image assigned
to carriers $X,Y$ is the closure of $\phi(X\times Y)$.

\begin{lemma}[Product testing]\label{prof:product-testing}
Take two pieces, or the pooled pieces described above, in a common
window with the same limiting profile. On a fixed plateau, fix in
each set an assignment of every point to one carrier of its original
piece that contains it. These assignments precede the tests below
and are retained under later restrictions. Let $X,Y$ denote the
assigned carriers of a pair, and fix an interior testing time $b$.
A polynomial of degree $O(T_s(b))$ captures only
$o(n_xn_y)$ pairs at which it does not vanish identically on
$X\times Y$. The same holds for a number of rational tests bounded
independently of the sequence, provided their denominators are
nonzero at the pairs concerned and their numerator polynomials have
degree $O(T_s(b))$ on the chosen charts, with a common implied
constant.

Consequently an equation system containing the labels of a dense
graph contains the complete assigned component images at all but
a negligible fraction of its pairs, whenever the label map has
bounded degree on these charts and the cleared numerator polynomials
of all pulled-back equations have degree $O(T_s(b))$ with a common
implied constant. The denominators used in these pullbacks must be
nonzero at the actual pairs concerned. For any collection of rational
rank tests whose size is bounded independently of the sequence and
which are generically nonzero on the relevant assigned images, the
tests are nonzero at all but $o(n_xn_y)$ corresponding actual pairs,
provided their denominators are nonzero there and their cleared
pullback numerators obey the same degree bound.
\end{lemma}

\begin{proof}
For an actual source point $p$, apply the individual-cut conclusion
of Proposition~\ref{prof:regularity} on every target carrier $Y$
on which the specialization is nonidentically zero.
For the other case, fix $Y$.  If the polynomial is not identically
zero on $X\times Y$, choose a generic $y\in Y$ for which its
specialization is nonidentically zero on each such $X$.
The same individual-cut conclusion on the source bounds points $p$ where the polynomial vanishes
identically on $Y$.  Average this bound with the target masses of
the $Y$'s. The fixed assigned sets are disjoint within each factor,
so these weighted sums are unambiguous even when carriers meet.
For pooled pieces, retain their original mass accounting.
All bounds are uniform over the polynomial and the
pieces by their sequence formulation.

For a system of equations, pull back each equation and clear its
chart denominator separately. The resulting numerator polynomials
vanish at the actual system points and have the common degree bound.
On every component product carrying a pair on the chart whose image
is not contained in the system, at least one numerator is
nonidentically zero. A generic linear combination of the numerators
is therefore nonidentically zero on every such product
simultaneously. It has the same degree bound: a finite basis of
their span suffices, regardless of the number of system equations.
Apply the first assertion to this combination.

For rank tests use their cleared pullback numerators. When a rank
condition is expressed by minors whose cleared numerators obey the
common degree bound, a generic linear combination of those numerators
is nonidentically zero on every relevant component product where the
rank condition holds generically, and its zero set contains every
rank failure. A bounded number of the resulting
tests has a union of failure sets of size $o(n_xn_y)$. A slightly
later interior time absorbs the uniform constant in the degree
bound. This proves the last assertion without assuming that sample
points are generic on their carriers.
\end{proof}

\subsection{Selecting the two windows}

\begin{proposition}[Selection of profile pairs]\label{prof:pair-selection}
For every $p\in P$, let $C_p$ be a finite family of irreducible
projective curves in $\PP^{d-1}$, each spanning a projective subspace
of dimension at most two, with total degree at most $E=\delta N/A$.
Suppose that $1-o(1)$ of
the unordered pairs $\{p,q\}$ satisfy $\pi_p(q)\in C_p$ or
$\pi_q(p)\in C_q$.  After the global terminal deletion described
below, one of the following configurations occurs.
\begin{enumerate}
\item A core of positive relative size is covered.
\item Two sequences of non-core pieces have a common dyadic width
$s$, the same limiting profile, and an ordered product covered in
one orientation with positive relative density.
\item Two sequences of non-core pieces have widths with ratio
tending to infinity, both $o(\log B)$, the same limiting profile,
and such a positively covered ordered product.
\end{enumerate}
In the common outermost window $s=\log B$, the two factors may
instead be the same pool of pieces in a shrinking profile bin.
That pool has size $n\ge N\sqrt\delta$, is the union of
$\exp(o(s))$ pieces, and retains all the preceding degree, mass,
regularity, and product-testing conclusions to exponential
accuracy.  All selected incidences may be required to lie on
physical lines containing at most $B^{o(1)}$ sample points.

The global terminal deletion means the following.  Whenever a
target point was assigned to its carrier curve after the
$(d-1)$st cut, its projection from the center is not a curve
contained in the center's directional family.
\end{proposition}

\begin{proof}
Assign each point of every non-core piece to one of its carrier
curves after the $(d-1)$st cut.  Omit assignments to nonreal lines.
Such a line contains at most one real point, so the omitted
fraction in a piece is at most
\[
 \frac{e_{d-1}\exp(o(s))}{n}
 =D_d^{-1}\exp(o(s))=o(1),
\]
using~\eqref{prof:balance}.  This is uniform over pieces.  Every
point is assigned at most once; pool repeated curves.  Apply
Proposition~\ref{geo:terminal} once to this global assignment.  It removes
only $o(N^2)$ ordered incidences, and hence leaves the required
almost-complete unoriented coverage.  Assignments retain their
original meaning if pieces are later pooled.

If the core has positive relative size, the cone over a
directional family has degree $O(E)$ and dimension at most two.
The core's isolation equations, restricted to its components,
and B\'ezout bound its intersections with the core by $O(EL^2)$
per center.  This is $o(N)$, so neither orientation can cover a
positive fraction of the core pairs.  This resolves the first
alternative.  Otherwise discard the core.

Here are details of the remaining selection, including the order
of the limits.  First fix a precision $1/h$, and partition the
compact space of normalized drop lists and values of
$s/\log B\in[0,1]$ into finitely many bins of that diameter.
Also bin the dyadic index modulo an integer tending to infinity
with $h$.  Two independent random retained points fall in one
common bin with probability bounded below by a positive constant
depending only on $h$.  Thus the relative uncovered probability
conditional on this agreement tends to zero for fixed $h$.
Select a pair of pieces with small relative uncovered
probability, and pass to an orientation containing at least half
the covered pairs.  Now let $h$ tend to infinity sufficiently
slowly with the sample index.  This standard diagonal choice
retains conditional coverage while making the profiles and
normalized widths converge to the same limits.

The selected dyadic indices are either equal on a subsequence,
or their difference tends to infinity, by their congruence
condition.  In the latter case the width ratio tends to infinity.
Their normalized widths have the same limit, so that limit must
be zero.  This proves the two window alternatives.

There is a useful refinement at the outermost window.  If its
total mass tends to zero, discard it before the selection.  If
its mass is bounded below on a subsequence, work only there.
For each fixed precision choose a largest profile bin and use
its pool as both factors.  It has a fixed positive mass at that
precision, and its relative uncovered fraction tends to zero.
In the diagonal choice increase precision slowly enough that
its mass is at least $N\sqrt\delta$ and
$N\exp(-o(\log B))$.  This is possible because $\delta\to0$.
Each constituent piece has size $N\exp(-o(s))$ uniformly, so
there are only $\exp(o(s))$ of them.  Their profiles differ by
$o(1)$, and the conclusions for pooled families follow from
Proposition~\ref{prof:regularity} and
Lemma~\ref{prof:product-testing}.

Finally, Szemer\'edi--Trotter, with a generic planar projection,
makes the fraction of ordered pairs on lines of occupancy above
a threshold $H$ arbitrarily small as $H\to\infty$: summing its
dyadic rich-line estimate gives $O(N^2/H+N h_{\max})$, and the
proper-flat bound gives $h_{\max}=o(N)$.  At each fixed selection
precision first choose $H$ large enough for its conditioning
probability.  Increase the precision sufficiently slowly that
$H\le B^{o(1)}$ and all the preceding conditional errors still
vanish.  This is compatible with the outer-pool mass condition.
All constants and precisions are fixed before taking each
asymptotic limit; the final diagonal sequence merely combines
these finitely specified requirements.
\end{proof}
